\documentclass{report}
\usepackage[utf8]{inputenc}
\usepackage{tikz}
\usepackage{standalone}
\usepackage{amsmath}
\usepackage{amsfonts}
\usepackage{graphicx, color}
\usepackage[a4paper,margin=2cm]{geometry}

\newcommand{\red}[1]{{\color{red}{#1}}}


\newcommand{\Norm}{\mathcal{N}}
\newcommand{\Loss}{\mathcal{L}}
\newcommand{\R}{\mathbb{R}}
\newcommand{\bt}{\mathbf{t}}
\newcommand{\bU}{\mathbb{U}}
\newcommand{\bu}{\mathbf{u}}
\newcommand{\bw}{\mathbf{w}}
\newcommand{\bX}{\mathbf{X}}
\newcommand{\bx}{\mathbf{x}}
\newcommand{\by}{\mathbf{y}}
\newcommand{\bZ}{\mathbf{Z}}
\newcommand{\bz}{\mathbf{z}}

\newcommand{\eq}{=}
\newcommand{\parfrac}[2]{\frac{\partial #1}{\partial#2}}


\title{Causal Effect Inference using Normalizing Flows}
\author{Micha de Groot}


\begin{document}

%%%%%%%%%%%%%%%%%%%%%%%%%%%%%%%%%%%%%%%%%%%%%%%%%%%%%%%%%%%%%%%%%%%%%%%%%%%%%%%%
\begin{titlepage}

\newcommand{\HRule}{\rule{\linewidth}{0.5mm}} % Defines a new command for the horizontal lines, change thickness here
\center % Center everything on the page
 
%----------------------------------------------------------------------------------------
%	HEADING SECTIONS
%----------------------------------------------------------------------------------------

\includegraphics[width=\linewidth]{uvaENG}\\[2.5cm]
\textsc{\Large MSc Artificial Intelligence}\\[0.2cm]
% \textsc{\normalsize Track: \red{track}}\\[1.0cm] % track
\textsc{\Large Master Thesis}\\[0.5cm] 

%----------------------------------------------------------------------------------------
%	TITLE SECTION
%----------------------------------------------------------------------------------------

\HRule \\[0.4cm]
{ \huge \bfseries Causal Effect Inference\\ with Normalising Flows}\\[0.4cm] % Title of your document
\HRule \\[0.5cm]
 
%----------------------------------------------------------------------------------------
%	AUTHOR SECTION
%----------------------------------------------------------------------------------------

by\\[0.2cm]
\textsc{\Large Micha de Groot}\\[0.2cm] %you name
10434410\\[1cm]


%----------------------------------------------------------------------------------------
%	DATE SECTION
%----------------------------------------------------------------------------------------

{\Large \today}\\[1cm] % Date, change the \today to a set date if you want to be precise

48 EC\\ %
October 2019 - June 2020\\[1cm]%

%----------------------------------------------------------------------------------------
%	COMMITTEE SECTION
%----------------------------------------------------------------------------------------
\begin{minipage}[t]{0.4\textwidth}
\begin{flushleft} \large
\emph{Supervisor:} \\
Dr. Efstratios Gavves% Supervisor's Name
\end{flushleft}
\end{minipage}
~
\begin{minipage}[t]{0.4\textwidth}
\begin{flushright} \large
\emph{Assessor:} \\
\red{Dr A  \textsc{Person}}\\
\end{flushright}
\end{minipage}\\[2cm]

%----------------------------------------------------------------------------------------
%	LOGO SECTION
%----------------------------------------------------------------------------------------

\framebox{\rule{0pt}{2.5cm}\rule{2.5cm}{0pt}}\\[0.5cm]
%\includegraphics[width=2.5cm]{figure}\\ % Include a department/university logo - this will require the graphicx package
\textsc{\large Instituut voor Informatica}\\[1.0cm] % 
 
%----------------------------------------------------------------------------------------

\vfill % Fill the rest of the page with whitespace

\end{titlepage}


\tableofcontents


%%%%%%%%%%%%%%%%%%%%%%%%%%%%%%%%%%%%%%%%%%%%%%%%%%%%%%%%%%%%%%%%%%%%%%%%%%%%%%%%
\chapter{Introduction}
Measuring the causal effect of an action on certain parts of the world and the people in it is a question that form the cornerstone of most sciences\footnote{\label{note:citation}Citation needed}. Most disciplines of science try to isolate and quantify such effects, to allow predictions of the future. One of the problems faced in practice is the difficulty in isolating the causal effect of an action from the influence of any background variables. 

In recent years the work by Pearl et al.\cite{pearl2009causal} \cite{pearl1995causal} has yielded a framework in which these causal effects can be modelled in terms of probability densities and in which it is theoretically possible to isolate a direct causal effect of a variable $\bt$ on an outcome variable $\by$ if there are one ore more confounding variables $\bZ$. Classically this is done through the use of randomised trials or other methods to isolate what is to be measured physically form any interference and by that preventing the existence of any possible confounder $\bZ$. Unfortunately this is not always possible and for those cases we shift our approach to the so-called \textit{do}-calculus. 

In cases where it is not possible to directly use equation \ref{equation:do_operation}, we use a proxy variable $\bx$ instead. The work of Louizos et al. \cite{louizos2017causal} has show that this can to some extend be done by using VAEs \cite{kingma2013auto}. This approach had the problem that the variational lower bound did not become as high as theoretically possible. To circumvent that problem we propose an alternative approach by using Normalising Flows \cite{rezende2016variational} to model the posterior distribution of $\bZ$ and use that to calculate the Average Treatment Effect.


%%%%%%%%%%%%%%%%%%%%%%%%%%%%%%%%%%%%%%%%%%%%%%%%%%%%%%%%%%%%%%%%%%%%%%%%%%%%%%%%
\chapter{Causal effect inference}
In the discipline of causal inference there are several questions that are commonly of interest. In this context we assume the framework of reasoning about causality as described by Pearl et al. \cite{pearl2009causal}, where relations between events are modelled as a Directed Acyclic Graph (DAG), and the state of each event is defined as a function of all its parents in said graph. What is interesting to know then are of course the structure of graphs and the form and parameters of each connection function. If the true graphs and functions would be know we would be able to precisely predict all causal effects. Of course such a bold claim follows by the observation that the search space of potential DAGs grows exponentially with the number of vertices, growing to $29281$ possible graphs when there are five vertices \cite{robinson1977counting}. Some methods have been developed to address this problem, under some assumptions, but that goes beyond the scope of this research.

In this research we assume a given structure of the DAG and focus on finding the relation between the random variables in the graph. Specifically the effect of one variable, called the treatment and denoted with $\bt$, on one other variable, called the outcome and denoted with $\by$. 

\section{The \textit{do}-calculus and its use in causal inference}
The rules of \textit{do}-calculus allow us to quantify the effect of $\bt$ on $\by$ when there are one or more confounding variables, which is usually the case. This only requires the probability distributions and the structure of the causal graph defining the relation between all variables. As we assume the graph structure to be known we only need a correct factorisation of the joint distribution of all variables and we are practically done. An example is drawn in Figure \ref{fig:graph_observed_confounder_and_latent_with_proxy}. In this graph we can see what the confounding variables are that cause the two variables who's effect we want to measure, and correct for that by using the famous \textit{do}-operator equation:

\begin{figure}
    \centering
    \includestandalone{Figures/observed_confounder}
    \hspace{2cm}
    \includestandalone{Figures/causal_graph_one_proxy_one_confounder}
    \caption{Two causal Bayesian graphs, where all observed variables are coloured grey, unobserved variables are white and the causal relations are represented as arrows. The left hand side models an observed confounder and the right hand side a latent confounder with a proxy variable.}
    \label{fig:graph_observed_confounder_and_latent_with_proxy}
\end{figure}


\begin{equation}\label{equation:do_operation}
   p(\by | do(\bt)) = \int_{pa_\bt} p(\by | \bt, pa_{\bt}) p(pa_{\bt}) \text{d} pa_\bt
\end{equation}
The set $\text{pa}_\bt$ denotes the set of ancestor nodes from $\bt$ that satisfy the backdoor-criterion\footnotemark[\ref{note:citation}], which is only $\bZ$ in the case of left hand side graph in Figure \ref{fig:graph_observed_confounder_and_latent_with_proxy}. The problem is that in most real world scenarios we don't know what $\bZ$ is exactly and how these probabilities are shaped: it is a \textit{latent} confounder\footnotemark[\ref{note:citation}]. Most scientific disciplines that implicitly use this framework, such a double-blind medical trials, circumvent this problem by eliminating all possible latent confounders by making the cause of $\bt$ independent of everything else: a random assignment.

But in general this is not possible, and instead the principle of proxy variables is introduced. These proxy variables, denoted with $\bX$, can be anything that can be measured and from which we can assume that it is directly caused by $\bZ$. The right hand side of Figure \ref{fig:graph_observed_confounder_and_latent_with_proxy} represent such a situation. The underlying principle of proxy variables is that it allows us to correct for the latent confounder in an approximate way. There has been extensive research in this area in recent years but most methods either assume that the latent confounder is categorical or 
% \cite{kuroki2014measurement} \cite{miao2018identifying}


%%%%%%%%%%%%%%%%%%%%%%%%%%%%%%%%%%%%%%%%%%%%%%%%%%%%%%%%%%%%%%%%%%%%%%%%%%%%%%%%
\chapter{Generative modelling and variational inference}
\section{Normalising Flows}
The research in generative models has yielded a model type called Normalising Flows, first thought of by Tabak and Turner \cite{tabak2013family} and later popularised by Rezende and Mohamed \cite{rezende2016variational}. The approach of this class of models is to learn an invertible series of mappings from a prior distribution of a simple form to the data likelihood. By using the change of variable rule, in Equation \ref{equation:change_of_variables}, it is guaranteed that before and after the transformations we have a valid probability distribution. Through the use of the inverse of these mappings one can perform exact posterior inference. 

\begin{equation}\label{equation:change_of_variables}
    \bx = f(\bz) \qquad p(\bx) = p(\bz) \left|\text{det} \parfrac{f}{\bz} \right|^{-1}
\end{equation}
The reason we talk about a series of transformations is to split the the potentially complex mapping in Equation \ref{equation:change_of_variables} into smaller, simpler transformations. This results then in Equation \ref{equation:change_of_variables_log_chain}, given in log-space, as is conventional. Here we have $K$ functions $f_k$ mapping from latent variables $\bz_{k-1}$ to $\bz_k$ and ending with the mapping from $\bz_{K-1}$ to $\bx$.

\begin{equation}\label{equation:change_of_variables_log_chain}
    \ln p(\bx) = \ln p(\bz_0) - \sum\limits^K_{k=1}\ln \left| \text{det} \parfrac{f_k}{\bz_{k-1}} \right|
\end{equation}

To make this work in practice the (log)determinant of the Jacobian of each mapping $f_k$ has to computed efficiently. The most straightforward way to do this is to enforce that each mapping has a triangular Jacobian. This immediately solves the second practical criterion of having a tractable inverse of each $f_k$. Another more implicit requirement for our mappings is that they are parameterised functions on which we can use gradient descent methods to learn the parameters. Several implementations of Normalising Flows have been made so far, the simplest of which is the planar flow and radial flow

\subsection{Planar flow and radial flow}
In the original work of Rezende and Mohamed \cite{rezende2016variational}, two possible implementations were proposed. The first one is the planar flow, in which each mapping has the form:
\begin{equation}\label{equation:planar_flow}
    f(\bz) = \bz + \bu h(\bw^T\bz + b)
\end{equation}
where $\bw \in \mathbb{R}^D$, $\bu \in \mathbb{R}^D$ and $b \in \mathbb{R}$ are learnable parameters and $h(\cdot)$ is a smooth element-wise non-linearity with derivative $h'(\cdot)$. The log-likelihood under a series of such transformations is defined as:
\begin{equation}\label{equation:planar_flow_logdet}
    \ln p(\bx) = \ln p(\bz_0) - \sum\limits^K_{k=1} \ln \left|1 + \bu_k^T h'(\bw_k^T \bz_{k-1} + b) \right|
\end{equation}

Each transformation here can be seen as a layer in a neural network that consists of a skip connection and a single-node dense layer followed by an expansion back to the original number of dimensions. The downside of this is the limited transformative capabilities of each mapping in the flow.

\subsection{Real-valued Non-Volume Preserving transformations}
A type of Normalising Flow is the Real-valued Non-Volume Preserving transformations (real NVP) \cite{dinh2016density}. By using so-called coupling layers, this model type encompasses a more powerful type of Normalising Flows. Each transformation in this model consists of two coupling layers, where each coupling layers transforms one half of the current variable vector $\bz_k \in \mathbb{R}^D$ and keeps the other half fixed, done in the following way:
\begin{align}\label{equation:real_nvp_coupling}
    \bz_{k+1, 1:d} &= \bz_{k, 1:d} \\
    \bz_{k+1, d+1:D} &= \bz_{k, d+1:D} \odot \exp \left(s(\bz_{k, 1:d}) \right) + t(\bz_{k, 1:d})
\end{align}
where $s$ and $t$ are scale and translation functions respectively, from $\mathbb{R}^{d} \rightarrow \mathbb{R}^{D-d}$. By having one coupling layer transforming $\bz_{d+1:D}$ and the next layer transforming $\bz_{1:d}$ the whole variable is transformed. The Jacobian of one coupling layer is triangular:
\begin{equation}
    \parfrac{\bz_{k+1}}{\bz_k^T} = \begin{bmatrix}
    \mathbb{I}_d & 0\\
    \parfrac{\bz_{k+1, d+1:D}}{\bz_{k, 1:d}^T} & \text{diag}(\exp(s(\bz_{k, 1:d}))
    \end{bmatrix}
\end{equation}
which gives an easy to compute log determinant: $\sum\limits^d_{i=1} s(\bz_{k, 1:d})_i$. The log-determinant Jacobian  does not require us to compute a Jacobian or determinant of either $s$ or $t$. Computing the inverse of each coupling layer doesn't require the inverse of $s$ or $t$ either, as we only need to invert the multiplication and addition:
\begin{align}\label{equation:real_nvp_coupling_inverse}
    \bz_{k, 1:d} &= \bz_{k+1, 1:d} \\
    \bz_{k, d+1:D} &= (\bz_{k+1, d+1:D} - t(\bz_{k+1, 1:d})) \odot \exp \left(- s(\bz_{k+1, 1:d}) \right) 
\end{align}
The simplicity of these equations allow us to choose $s$ and $t$ arbitrarily complex, by choosing a deep neural network for example. The split of each vector $\bz_k$ into two halves can be done arbitrarily, not requiring the elements of the two halves to be consecutive elements. The pattern in from which the two halves are constructed have a variety op options. The authors of \cite{dinh2016density} suggest to  either use a checkerboard pattern, if the data consists of images, or to reshape the input to contain a multiple of the original number of channels and alternate between channels to which half of the split they belong. In all cases is is pertinent to transform values every other coupling layer.



%%%%%%%%%%%%%%%%%%%%%%%%%%%%%%%%%%%%%%%%%%%%%%%%%%%%%%%%%%%%%%%%%%%%%%%%%%%%%%%%

\chapter{Causal effect inference with generative models}
The idea of using a generative model in causal effect inference was proposed by Louizos et al. \cite{louizos2017causal}. They leveraged the power of VAEs to make it possible to infer latent confounders from data without any prior knowledge on the confounder itself.

\begin{figure}
    \centering
    \includestandalone[width=0.49\textwidth]{Figures/cevae_encoder_figure}
    \includestandalone[width=0.49\textwidth]{Figures/cevae_decoder_figure}
    \caption{The graphical model of the CEVAE with the encoder on the left hand side and the decoder on the right hand side. Grey nodes indicate sampling steps and white nodes indicate deterministic neural networks. In both the encoder and the decoder there is a clear distinction between the two values of $\bt$ for $\by$ and $\bz$ but in practice these are computed in one operation.}%Nog een ref naar formule oid?
    \label{fig:cevae_graphical_model}
\end{figure}

\section{Causal effect inference with Normalising Flows}
This can be potentially be improved upon by using Normalising Flows. Their ability to reproduce the true posterior over the latent confounder is an improvement on the CEVAE


%%%%%%%%%%%%%%%%%%%%%%%%%%%%%%%%%%%%%%%%%%%%%%%%%%%%%%%%%%%%%%%%%%%%%%%%%%%%%%%%
\bibliography{references.bib}
\bibliographystyle{apalike}


\end{document}